\documentclass[12pt, a4paper]{article}

\usepackage[T1]{fontenc}
\usepackage[utf8]{inputenc}
\usepackage[british]{babel}
\usepackage{mathtools}
\usepackage{amsthm}
\usepackage{libertine}
\usepackage[libertine]{newtxmath}
\usepackage[mathscr]{euscript}
\usepackage{listings}
\usepackage{enumitem}
\usepackage{tikz-cd}
\usetikzlibrary{decorations.markings}
\usepackage{float}
\usepackage{geometry}
 \geometry{
 a4paper,
 total={170mm,257mm},
 left=20mm,
 top=30mm,
 }
\usepackage{fancyhdr}
\usepackage{hyperref}
\urlstyle{same}
\usepackage[noabbrev]{cleveref}
\usepackage{xcolor}
\usepackage{tcolorbox}
\tcbuselibrary{listings}

% Colors to match VS Code Lean 4 Light Theme
\definecolor{keywordcolor}{rgb}{0.0, 0.0, 1.0}   % blue
\definecolor{tacticcolor}{rgb}{0.0, 0.0, 1.0}    % blue
\definecolor{commentcolor}{rgb}{0.0, 0.5, 0.0}   % green
\definecolor{symbolcolor}{rgb}{0.0, 0.0, 0.0}    % black
\definecolor{sortcolor}{rgb}{0.0, 0.0, 1.0}      % blue
\definecolor{attributecolor}{rgb}{0.0, 0.0, 0.0} % black
% Light gray background color
\definecolor{codebg}{rgb}{0.97, 0.97, 0.97}

\def\lstlanguagefiles{lstlean.tex}

\newtcblisting{leancode}{
  listing only,
  listing options={
    language=lean,
    emph={sorry},
    emphstyle=\color{red},
    basicstyle=\ttfamily\footnotesize
  },
  colback=codebg,
  colframe=codebg,
  boxrule=0pt,
  arc=3mm,
  auto outer arc,
  top=10pt,
  bottom=10pt,
  left=10pt,
  right=10pt,
  width=0.9\textwidth,
}

\theoremstyle{definition}
\newtheorem{innercustomthm}{Exercise}
\newenvironment{xca}[1] {\renewcommand\theinnercustomthm{#1}\innercustomthm}
  {\endinnercustomthm}

\pagestyle{fancy}
\fancyhf{}
\rhead{Lean Algebra Exercises{\renewcommand{\thefootnote}{\fnsymbol{footnote}}\footnotemark[2]}}
\chead{Group Homomorphism}
\lhead{Pedro N\'{u}\~{n}ez}
% \rfoot{Page \thepage}

\title{Lean Algebra Exercise 1.2}
\author{Pedro N\'{u}\~{n}ez}

\setcounter{tocdepth}{1}
\sloppy
\makeatletter
\hypersetup{
  pdfauthor={\@author},
  pdftitle={\@title},
  colorlinks,
  linkcolor=[rgb]{0.2,0.2,0.6},
  citecolor=[rgb]{0.2,0.2,0.6},
  urlcolor=[rgb]{0.2,0.2,0.6}}
\makeatother

\begin{document}

\begin{xca}{1.2}
  Let $G$ and $G'$ be groups with neutral elements $e \in G$ and $e' \in G'$.
  Let $f \colon G \to G'$ be a group homomorphism.
  Show that:
  \begin{enumerate}[label=(\alph*)]
    \item One has $f(e) = e'$.
    \item For all $a \in G$, one has $f(a^{-1}) = f(a)^{-1}$.
  \end{enumerate}
  \href{https://leanprover-community.github.io/}{Lean} is a computer program that one can use to verify mathematical proofs.
  Our colleague Johan Commelin is currently working with Peter Scholze to verify the proof of one of his theorems with Lean.
  You can try to \href{https://live.lean-lang.org/#codez=JYWwDg9gTgLgBAWQIYwBYBtgCMB0BBdAcwFMsokcBxKCAVzBwCEkBnYAYwChOBaHuAGLAoLeABNiAM2AA7YnCRx26VizjEZ7CGNmE4aeYRr04qCCHPQwqYCxAKAHsHM5Oy1XGp0wACXNwAb0pPAHI4AC44ABUATzBiAF84AG0vE0oAXRS0sFCsgApJCM84QCTCUIBKYoAFGlyAd1RiKGJOOFNzAH0QWnRiwAAiBTgsABo4IvzFACphqoBecaGZoqxuPmimmThRFHkkGTE4MBoAN3kDOAk2FsOWll6YVwNoYntQJBJOiEkvuUDgyhhSKxeJJVLGXKZbIQvJtOCFYrBcqAqrg7x%2BeySLKRIoARjgC3x4QWWBicJY0CgMTg6wAouQWPIKVAqQoDnAWmAVOx5MB4PU%2Bag4DE6FA4BB6ltjhBvgBCIA}{verify part (a) of this exercise with Lean}.
\end{xca}

\subsection*{Hints for Lean}

There are two ways to work with Lean:
\begin{itemize}
  \item Directly on your web browser:~\href{https://live.lean-lang.org}{live.lean-lang.org}.
  \item Locally on your computer:
    \begin{enumerate}
      \item First install Lean:~\href{https://lean-lang.org/install}{lean-lang.org/install}.
      \item Then create a new Lean project: \\
        \href{https://leanprover-community.github.io/install/project.html#creating-a-lean-project}{leanprover-community.github.io/install/project.html\#creating-a-lean-project}.
    \end{enumerate}
\end{itemize}
Once you are ready, you can start as follows:

\begin{leancode}
  import Mathlib.Algebra.Group.Basic

  -- First define a class encoding the group homomorphism axiom.
  class GroupHom {G G' : Type} [Group G] [Group G']
    (f : G → G') : Prop where
    hom_mul : ∀ a b, f (a * b) = f a * f b

  -- Then state and prove the desired result.
  theorem image_of_one {G G' : Type} [Group G] [Group G']
    (f : G → G') [GroupHom f] : f 1 = 1 := by
    sorry -- Erase sorry and replace it with your own proof!
\end{leancode}

Of course, this is not enough information to complete the task if you are using Lean for the first time.
The website \href{https://leanprover-community.github.io}{leanprover-community.github.io} contains many tutorials and learning resources.
If you want to jump right into it, you can copy-paste the contents of the file \href{https://github.com/pedro-nunez/lean-algebra-2021/blob/master/group-homomorphism/hints.lean}{hints.lean} available on \href{https://github.com/pedro-nunez/lean-algebra-2021/tree/master/group-homomorphism}{github.com/pedro-nunez/lean-algebra-2021/tree/master/group-homomorphism} into Lean's web editor \href{https://live.lean-lang.org}{live.lean-lang.org} and have a look at the commented example given there.
You can also find a \href{https://github.com/pedro-nunez/lean-algebra-2021/blob/master/group-homomorphism/solution.lean}{solution} on the repository linked above.

{\renewcommand{\thefootnote}{\fnsymbol{footnote}}\footnotetext[2]{Exercises designed for the lecture \href{https://home.mathematik.uni-freiburg.de/arithgeom/lehre/ws21/azt}{\emph{Algebra and Number Theory}}, for which I was a teaching assistant.
This was an introductory lecture on abstract algebra, polynomial rings and Galois theory, taught by Annette Huber-Klawitter at the University of Freiburg during the Winter Semester 2021/2022.}}

\end{document}
